% Source: https://artofproblemsolving.com/community/c3h1349381_what_is_the_probability_that_the_product_is_0
% Stable identity: aops-topic-1349381-q01
A product of two numbers is $0$ exactly when one of them is $0$ (all the other numbers in the set are nonzero).

(i) Unordered pairs. There are $\dbinom62=\dfrac{6\cdot5}{2}=15$ equally likely pairs. The pairs containing $0$ are $0$ together with any of the other $5$ numbers, so $5$ pairs. The probability is $\dfrac{5}{15}=\dfrac{1}{3}$.

(ii) Ordered pairs. There are $6\cdot5=30$ equally likely ordered choices (the second number must differ from the first). The number $0$ can be chosen first ($5$ choices for the second number) or second ($5$ choices for the first number), giving $10$ favourable ordered pairs, and the probability is $\dfrac{10}{30}=\dfrac{1}{3}$.

Both methods agree because each unordered pair corresponds to exactly $2$ ordered pairs, so both the total and the favourable count are doubled in (ii). The probability is $\boxed{\dfrac{1}{3}}$.
